Real-world urban counter-UAV deployments rarely proceed without disruption. Problem 4 embeds our frozen nominal $k=8$ schedule into a 60 s operational horizon disturbed at $t_d=5.0$ s by three classes of stochastic hazards—single-vehicle mechanical failure, sudden appearance of a fourth alert T4 (altitude 50 m, horizontal range $30$--$50$ m from $O$), and dispatch-center communication blackout—appearing either individually or in compound combinations. To evaluate recourse without introducing post-hoc scheduling bias, we build a deterministic discrete-event simulation engine \cite{stentz1994,vantendeloo2018} that replays every scenario from the identical frozen baseline (Algorithm~\ref{alg:p4-reschedule}).

\subsection{Disturbance Clock Semantics and Cached-Plan Sensitivity}
Rigorous event simulation requires unambiguous separation between three timekeeping concepts that frequently cause hidden discrepancies in operational studies: (i) the \emph{disturbance occurrence clock} ($t_d=5.0$ s), when a fault, new alert, or link drop begins; (ii) the \emph{absolute recovery clock} $t_{\mathrm{rec}}$ versus the \emph{outage duration} $\Delta t_{\mathrm{out}}=t_{\mathrm{rec}}-t_d$; and (iii) the \emph{onboard cached-instruction semantics} governing autonomous vehicle behavior while disconnected. Under the benchmark's primary clock formula, $t_{\mathrm{rec}}\in(5,30]$ s corresponds to an upper outage duration of $\Delta t_{\mathrm{out}}=25.0$ s; to resolve potential textual ambiguity with a 30 s outage-duration interpretation, our verification suite explicitly executes the $\Delta t_{\mathrm{out}}=30.0$ s ($t_{\mathrm{rec}}=35.0$ s) sensitivity branch alongside the primary suite, confirming that compound recovery shifts predictably by $+5.000$ s (from $34.826$ s to $39.826$ s under combined V2 fault, T4 alert, and blackout).

Similarly, onboard schedule caching proves decisive during pure communication loss. When vehicles cache the complete nominal schedule prior to link failure at $t_d=5.0$ s, the fleet completes all three alerts autonomously by $11.034$ s even under a 25 s blackout ($t_{\mathrm{rec}}=30.0$ s). By contrast, under an \emph{ongoing-only} non-cached instruction model where vehicles finish only their currently active leg and wait for link restoration before initiating reused slots, the latest paired observation start degrades sharply to $33.727$ s—a $+22.693$ s latency penalty that directly quantifies the engineering return on lightweight onboard plan caching.

\subsection{Lexicographic Evaluation Across All 73 Verified Disturbance Scenarios}
At the dispatch release clock $t_{\mathrm{rel}}$ (equal to $t_d=5.0$ s when communications remain active, or $t_{\mathrm{rec}}$ upon link restoration), alerts whose paired observation has already begun with two healthy vehicles are permanently locked ($\mathcal{A}_{\mathrm{lock}}$). Under \textbf{Keep Plan}, remaining vehicles continue executing cached nominal slots without central reassignment. Under \textbf{Reschedule}, any active vehicle not engaged in a locked observation may be reassigned to pending alerts ($\mathcal{A}_{\mathrm{pend}}$) after paying the mandatory gimbal cancellation-reset duration from its instantaneous pose $(\alpha_i(t_{\mathrm{rel}}),\beta_i(t_{\mathrm{rel}}))$ back to $(0^\circ,0^\circ)$ (Algorithm~\ref{alg:p4-reschedule}, lines 8--11). Both policies are compared lexicographically via Eq.~\eqref{eq:p4-score}: maximize completion fraction $\eta=k/N$ first (the share of alerts initiating valid paired observation within $[0,60]$ s), and break ties by minimizing the latest paired start $t_{\mathrm{last}}$ (censored at 60 s if any alert remains unstarted).

Table~\ref{tab:p4cat} and Figure~\ref{fig:p4} summarize paired outcomes across the complete suite of $\BCaseCount{}$ verified disturbance case IDs spanning all seven single-hazard and compound categories (comprising 28 archived benchmark cases and 45 structured cross-stress cases, including 11 overlapping parameter tuples retained transparently for audit completeness rather than extrapolated as independent samples over a continuous disturbance distribution). The operational advantage of dynamic rescheduling is decisive: Keep Plan averages only $\eta=0.604$ completion across the suite—failing completely on every new-alert scenario and on four of five single-vehicle faults—whereas Reschedule achieves $\eta=1.000$ (100\% alert coverage within the 60 s horizon) on all \BCaseCount{}/\BCaseCount{} scenarios. Lexicographically, Reschedule records \textbf{68 wins, 4 ties, and exactly 1 loss} against Keep Plan. The four ties occur in pure cached-plan outages (3 cases, where cached execution already finishes at $11.034$ s before link recovery) and one compound fault-outage case where both policies achieve identical timing after recovery. Only a single scenario—V5 single-vehicle failure at $t=5.0$ s—produces a slowdown under unconditional rescheduling, motivating the mechanistic dissection below.

\ReviewFigure{fig07_p4_resilience.pdf}{98mm}
{Paired completion fraction, censoring-aware timing contrast, and lexicographic score matrix across all \BCaseCount{} disturbance cases with event-driven rescheduling flow.}
{fig:p4}
